\chapter{Первый гейт ярусного родительского действия}

Предложенная архитектура различает интегрально нормированные компактные
координаты и \(L^2\)-нормированные распространяющиеся амплитуды. Это
различение уже имеет точное содержание в нейтринной суперсвязности:
нулевая форма \(\widehat e_0\) является нормированной волновой функцией, а
одна форма \(e_1\) имеет единичный период. Необходимо проверить, переносится
ли то же правило на заряженный лептонный сектор без второго скрытого
считывания меры.

\section{Алгебраический тест трёх чисел}

Доступные нормировки действительно воспроизводят три требуемые величины:
\begin{align}
 \rho_\tau^{\mathrm{raw}}
 &=\|1_{\mathbb{RP}^3}\|^2+
   \|1_{S^1}\|^2+
   \left\langle\|P_\perp n\|^2\right\rangle_{S^2}
 =\pi^2+2\pi+\frac23,\\
 \rho_\tau^{\mathrm{can}}
 &=1+1+\frac23=\frac83,\\
 \|\boldsymbol\Xi_\nu\|^2
 &=23\|\widehat e_0\|^2+\|e_1\|^2
 =23+\pi^{-1}.
\end{align}
Следовательно, числовой P1-гейт проходит как тождество норм. Однако это
ещё не означает, что все три значения являются выводами одного действия.

\section{Гейт типовой классификации}

Квантованный период имеет только нейтринная одна форма:
\begin{equation}
 \oint_\gamma e_1=1.
\end{equation}
Напротив, \(1_{\mathbb{RP}^3}\) и \(1_{S^1}\) являются постоянными
нуль-формами. Их амплитуды не имеют интегрального периода и не являются
вильсоновскими координатами только потому, что носитель компактен. При
канонической нормировке
\begin{equation}
 1_{\mathbb{RP}^3}\mapsto\pi^{-1}1_{\mathbb{RP}^3},
 \qquad
 1_{S^1}\mapsto(2\pi)^{-1/2}1_{S^1},
\end{equation}
и два объёмных вклада превращаются в \(1+1\).

Тем самым детерминированное правило
\begin{quote}
интегральный период \(\Rightarrow\) компактная нормировка;
обычная амплитуда нуль-формы \(\Rightarrow L^2\)-нормировка
\end{quote}
даёт одновременно \(23+\pi^{-1}\) и \(8/3\), но не даёт сырой
лептонный seed. Правило, сохраняющее ненормированные постоянные функции,
даёт seed, но не объясняет, почему те же направления в вершине должны
читаться канонически.

\begin{nogocriterion}[Один ярус --- два лептонных считывания]
Одна фиксированная типовая классификация не может назначить двум
постоянным лептонным нуль-формам одновременно сырые нормы
\((\pi^2,2\pi)\) и канонические нормы \((1,1)\). Для получения обоих
результатов требуется второе отображение считывания или перенормировка
вершины. Пока это отображение не выведено из симметрии, граничного действия
или квантовой меры, оно является тем же скрытым относительным весом, который
запрещён двухсекторным гейтом.
\end{nogocriterion}

Число \(8/3\) в этом тесте не является новым наблюдаемым сектором. Оно
остаётся контрольным результатом канонической нормировки, которая ранее
уничтожила сырой seed. Воспроизведение двух альтернативных координат одного
направления не считается двумя физическими предсказаниями.

\section{Статус архитектур}

Стратифицированная суперсвязность сохраняет нейтринный результат, но в
текущем виде не выводит лептонный seed. Слой БФ/граничного ротора может
оставаться независимым кандидатом на петлевой функционал, однако он не
устраняет древесную неоднозначность без явного отображения между фоновыми
модулями и нормированными лептонными вершинами. Поэтому гибрид A+B не
закрыт отрицательно как будущая модель, но ещё не является одним
родительским действием.

\begin{proposition}[Вердикт P1]
P1 проходит алгебраический гейт трёх чисел и проваливает гейт единственной
меры и типовой классификации:
\begin{equation}
 N_{\mathrm{predictive\ sectors}}=1<2.
\end{equation}
Единственным прошедшим нормировочно-чувствительным сектором остаётся
нейтринная коллективная жёсткость. Нового закрытого физического
предсказания не возникает.
\end{proposition}

\begin{protocol}[Условие переоткрытия A+B]
Необходимо выписать локальное или граничное действие, которое до сравнения
с \(m_\tau\):
\begin{enumerate}
 \item выводит отображение фоновых компактных модулей в канонически
 нормированную лептонную вершину;
 \item совместно преобразует stiffness и coupling при смене координат;
 \item фиксирует, почему древесная норма сохраняет
 \(\pi^2+2\pi+2/3\), а петлевая вершина получает единственный заранее
 определённый вес;
 \item тем же правилом проходит независимый электромагнитный или
 роторный гейт.
\end{enumerate}
До этого P2--P4 являются тестами отдельных слоёв, а не тестами одного
родительского действия.
\end{protocol}

Машинный протокол сохранён в
\texttt{s2t\_tiered\_parent\_action\_p1\_results.json}.